\section{GQA (Grouped Query Attention) explicado en detalle}

\subsection{Motivación y mecanismo de GQA}

\begin{importantbox}{Idea central de GQA}
En la MHA estándar, cada cabeza de Query corresponde a una cabeza de KV independiente. GQA hace que varias cabezas de Query \textbf{compartan} un mismo grupo de cabezas de KV, con lo cual se reduce considerablemente el costo de almacenamiento de la KV Cache.
\end{importantbox}

Tomando como ejemplo Qwen3-30B-A3B:
\begin{itemize}
  \item 32 cabezas de Query, 4 grupos de KV
  \item Cada $32 / 4 = 8$ cabezas de Query comparten un grupo de $K, V$
  \item Cabezas de Query 0--7 $\to$ grupo de KV 0; cabezas de Query 8--15 $\to$ grupo de KV 1; y así sucesivamente
\end{itemize}

\subsection{$d_{\text{model}} \neq n_{\text{heads}} \times d_{\text{head}}$}

\begin{warningbox}{Las dimensiones no siempre coinciden}
En LLaMA 2 o GPT-2, $d_{\text{model}} = n_{\text{heads}} \times d_{\text{head}}$. Pero en el modelo Dense de Qwen3, $d_{\text{model}} = 5120$, mientras que $n_{\text{heads}} \times d_{\text{head}} = 64 \times 128 = 8192 \neq 5120$. Esto significa que $W_Q$ ya no es una matriz cuadrada, sino una matriz no cuadrada de $5120 \times 8192$.
\end{warningbox}

\subsection{Resumen de la sección}

GQA comprime la KV Cache haciendo que varias cabezas de Query compartan grupos de KV, y es un estándar en los LLM predominantes actuales. Qwen3 va más allá y rompe la restricción de alineación entre $d_{\text{model}}$ y la dimensión de las cabezas, usando matrices de proyección no cuadradas.

